% =====================================================================
% 09_conclusion.tex — conclusion générale
% =====================================================================
\chapter{Conclusion}
\label{chap:conclusion}

Le projet réalise une base bancaire complète sur Oracle 21c XE, en combinant
le modèle relationnel et le modèle objet. Il couvre les trois volets du sujet :
les clients, conseillers et agences ; les comptes courants, d'épargne et à
terme ; les opérations sur comptes courants. Il ajoute les mécanismes
programmables demandés : fonctions, procédures, méthodes, triggers, type
récursif, collection multi-niveaux et vue objet.

La partie A établit un schéma contraint et rejouable. Les types métier sont
créés dans l'ordre des dépendances, les comptes reposent sur une super-classe
non instanciable, les téléphones utilisent un \code{VARRAY} et les signataires
une \emph{nested table}. Les contraintes ne sont pas décoratives : elles
encodent le statut, le découvert, le taux, la durée et le titulaire. La
contrainte différée de \code{cpt\_courant} prépare B5b sans affaiblir la règle
métier.

La partie B montre que le schéma supporte des manipulations réelles. Le client
99 est créé avec un seul téléphone, puis \code{change\_Portable} étend le
\code{VARRAY}. CC99 est inséré sans titulaire dans une transaction, puis mis à
jour avant le \code{COMMIT}. La date d'opération est prise par défaut, un
signataire est ajouté dans la \emph{nested table} et \code{drop\_client}
supprime toutes les dépendances. Ces scénarios vérifient à la fois les
contraintes et les limites du modèle.

La partie C démontre que le modèle objet reste interrogeable. Les \code{REF} se
déréférencent, les collections se déplient avec \code{TABLE(...)}, les
agrégats et les fonctions de classement restent disponibles. Les résultats
montrent un compte courant à découvert, un top trois des soldes, des échéances
proches et une opération de plus d'un million de dinars. La chronologie est
essentielle : ces résultats sont ceux d'avant les tests destructifs de la
partie D.

La partie D transforme le schéma en système actif. Les sous-programmes
encapsulent les soldes, les ajouts de clients et les virements. Les blocs
anonymes illustrent les curseurs et les collections. Les méthodes ajoutent le
polymorphisme de \code{fraisTenue}, le calcul des intérêts et la clôture de
compte. Les triggers garantissent le statut, les droits, le plafond, le
découvert, la mise à jour du solde et la traçabilité. Le trigger
\code{COMPOUND} résout proprement le problème de la table mutante.

La partie E complète le dispositif avec trois mécanismes avancés. Le type
récursif permet de naviguer de parrain en grand-parrain. La collection
multi-niveaux stocke un relevé annuel dans une seule colonne et reste
interrogeable par double \code{TABLE}. La vue objet agrège les conseillers et
les comptes d'une agence, et sa méthode \code{mutation} ferme puis rouvre une
affectation en conservant l'historique.

Le projet atteint ses objectifs techniques. Le script complet a été rejoué de
zéro, sans erreur de compilation, avec zéro objet invalide. Les seules erreurs
Oracle présentes dans les sorties sont des tests d'exception volontaires. Les
contraintes, les triggers et les méthodes se répondent : les règles métier ne
sont pas seulement documentées, elles sont exécutées par la base.

Le projet a aussi des limites, assumées et documentées. Les opérations portent
uniquement sur les comptes courants, car \code{operation\_type.ref\_compte}
est un \code{REF cptCourant\_type}. Le virement relie deux comptes courants.
Le responsable d'agence n'est pas remplacé automatiquement lors d'une mutation.
Il n'y a ni interface graphique, ni tests de concurrence, ni séparation fine
des privilèges. Ces limites ne sont pas des oublis : elles délimitent le
périmètre du mini-projet et ouvrent les perspectives présentées au chapitre
suivant.

Sur le plan méthodologique, le projet a suivi une démarche incrémentale. Chaque
partie a été exécutée et relue avant la suivante. Les comportements Oracle ont
été testés sur l'installation réelle, ce qui a évité plusieurs impasses :
\code{CREATE TYPE} direct, \code{ORA-22979} sur les \code{REF} d'OID primaire,
\code{ORA-04091} sur la table mutante et \code{ORA-02290} sur les dates de
mutation. Les hypothèses de cadrage ont été confrontées au comportement de la
base, puis corrigées lorsque nécessaire.

Le résultat est un livrable reproductible : un script maître, des scripts
séparés par partie, un jeu de données, un guide de soutenance, une liste de
captures et un rapport. La base peut être remise dans son état initial, puis
rejouée intégralement. Cette propriété est importante pour l'évaluation et
pour la démonstration : elle garantit que les résultats présentés ne dépendent
pas d'un état caché.
